% Options for packages loaded elsewhere
\PassOptionsToPackage{unicode}{hyperref}
\PassOptionsToPackage{hyphens}{url}
\documentclass[
]{article}
\usepackage{xcolor}
\usepackage{amsmath,amssymb}
\setcounter{secnumdepth}{5}
\usepackage{iftex}
\ifPDFTeX
  \usepackage[T1]{fontenc}
  \usepackage[utf8]{inputenc}
  \usepackage{textcomp} % provide euro and other symbols
\else % if luatex or xetex
\fi
\usepackage{lmodern}
\ifPDFTeX\else
  % xetex/luatex font selection
\fi
% Use upquote if available, for straight quotes in verbatim environments
\IfFileExists{microtype.sty}{% use microtype if available
  \usepackage[]{microtype}
  \UseMicrotypeSet[protrusion]{basicmath} % disable protrusion for tt fonts
}{}
\makeatletter
\@ifundefined{KOMAClassName}{% if non-KOMA class
  \IfFileExists{parskip.sty}{%
    \usepackage{parskip}
  }{% else
    \setlength{\parindent}{0pt}
    \setlength{\parskip}{6pt plus 2pt minus 1pt}}
}{% if KOMA class
  \KOMAoptions{parskip=half}}
\makeatother
\usepackage{listings}
\newcommand{\passthrough}[1]{#1}
\lstset{defaultdialect=[5.3]Lua}
\lstset{defaultdialect=[x86masm]Assembler}
\usepackage{longtable,booktabs,array}
\newcounter{none} % for unnumbered tables
\usepackage{calc} % for calculating minipage widths
% Correct order of tables after \paragraph or \subparagraph
\usepackage{etoolbox}
\makeatletter
\patchcmd\longtable{\par}{\if@noskipsec\mbox{}\fi\par}{}{}
\makeatother
\ifLuaTeX
\else
\fi
\ifLuaTeX
\fi
\setlength{\emergencystretch}{3em} % prevent overfull lines
\providecommand{\tightlist}{%
  \setlength{\itemsep}{0pt}\setlength{\parskip}{0pt}}
% ═══ 由 环境/pdf/latex/md到LaTeX.py 生成 ═══
% 中文字体：思源宋体（子集，SIL OFL 1.1）；拉丁/数学：Latin Modern（随 TeX Live 分发）
% 编译：xelatex（本仓库自带 wasm 版 XeTeX 引擎，见 环境/pdf/latex/编译LaTeX.mjs）
\usepackage[T1]{fontenc}      % 否则 ASCII 的 < > | 在 OT1 下排成 ¡ ¿ —
\usepackage{lmodern}          % 必需：数学的物理字模（否则 xdvipdfmx 拒绝出 PDF）
\usepackage{amsmath,amssymb,mathtools}
\usepackage{booktabs,longtable,array,multirow}
\usepackage{graphicx}
\usepackage{xcolor}
\usepackage{listings}
% 不用 hyperref：新版 hyperref 强制依赖 bookmark.sty，而 bundle 里没有。
% 链接命令用降级定义，保证正文里的 \href / \url 不会报错。
\providecommand{\href}[2]{#2}
\providecommand{\url}[1]{\texttt{#1}}
\providecommand{\hypertarget}[2]{#2}
\providecommand{\hyperlink}[2]{#2}
\providecommand{\urlstyle}[1]{}
\providecommand{\texorpdfstring}[2]{#1}   % hyperref 的命令，剥离后仍被模板/标题用到
\providecommand{\phantomsection}{}
\providecommand{\pdfstringdef}[2]{}
\providecommand{\hypersetup}[1]{}

% ── 三套字体：中文 / 符号 / emoji（按字符 cmap 自动选择）──
\font\zhfont="[NotoSerifSC-Regular.ttf]:script=hani" at 10.5pt
\font\symfont="[DejaVuSans.ttf]" at 10.5pt
\font\emofont="[NotoEmoji-Regular.ttf]" at 10.5pt
\font\zhmonofont="[NotoSansSC-Regular.ttf]" at 9pt
% 用法：\zh{中文} / \sym{→} / \emo{🦙}
% 注意：必须**带参数**。写成 \def\zh{{\zhfont}} 是错的——字体赋值只在小群里有效，
% 那个小群随 \zh 立刻结束，于是中文仍然用拉丁字体排，报 Missing character。
\def\zh#1{{\zhfont #1}}
\def\sym#1{{\symfont #1}}
\def\emo#1{{\emofont #1}}
% 含汉字的行内代码：listings 的 \lstinline 遇到汉字必炸（`\lst@arg ->判`
% 未定义控制字，catcode 设 12 也无效），所以直接自己排。
\def\codezh#1{{\zhmonofont #1}}

\def\zhglue{\hskip0pt plus .06em minus .01em}
\linespread{1.12}
\setlength{\parindent}{0pt}
\setlength{\parskip}{0.45em}

% ── 代码块：等宽 + 中文字体（listings 默认字体不含汉字）──
\lstset{
  basicstyle=\zhmonofont,
  breaklines=true,
  breakatwhitespace=false,
  columns=fullflexible,
  keepspaces=true,
  showstringspaces=false,
  frame=single,
  framesep=4pt,
  rulecolor=\color{gray!45},
  backgroundcolor=\color{gray!7},
  xleftmargin=6pt,
  extendedchars=true,
  tabsize=2,
  % 代码块里的 emoji：NotoSansSC 没有这些字形（listings 不做字符级回退）。
  % 用 escapeinside 开口子，转换器把 emoji 逐字包成 (*@{\emo{⭐}}@*)。
  % 试过 literate={⭐}{{\emofont ⭐}}1，会把 ASCII 字母也吞掉并报
  % `Improper alphabetic constant`，不可用。
  escapeinside={(*@}{@*)},
}
\renewcommand{\lstlistingname}{\zh{代\zhglue 码}}

% ── 表格线宽（booktabs 风格）──
\renewcommand{\arraystretch}{1.25}

% ── 标题样式 ──
\usepackage{titlesec}
\titleformat{\section}{\Large\bfseries}{\thesection}{0.6em}{}
\titlespacing*{\section}{0pt}{1.1em}{0.5em}
\urlstyle{same}
\hypersetup{
  pdftitle={\zh{课\zhglue 题}04-0.1B\zh{端\zhglue 到\zhglue 端}},
  pdflang={zh-CN},
  hidelinks,
  pdfcreator={LaTeX via pandoc}}

\title{\zh{课\zhglue 题}04-0.1B\zh{端\zhglue 到\zhglue 端}}
\author{}
\date{}
\begin{document}
\maketitle

{
\setcounter{tocdepth}{2}
\tableofcontents
}
\section{\texorpdfstring{\zh{课\zhglue 题}04 \zh{开\zhglue 题\zhglue 简\zhglue 报} \sym{·} 0.1B \zh{端\zhglue 到\zhglue 端\zhglue 预\zhglue 训\zhglue 练\zhglue （}\textbf{\zh{首\zhglue 次\zhglue 上\zhglue 机}}\zh{）}}{\zh{课\zhglue 题}04 \zh{开\zhglue 题\zhglue 简\zhglue 报} \sym{·} 0.1B \zh{端\zhglue 到\zhglue 端\zhglue 预\zhglue 训\zhglue 练\zhglue （\zhglue 首\zhglue 次\zhglue 上\zhglue 机\zhglue ）}}}\label{ux8bfeux989804-ux5f00ux9898ux7b80ux62a5-0.1b-ux7aefux5230ux7aefux9884ux8badux7ec3ux9996ux6b21ux4e0aux673a}

\begin{quote}
\zh{模\zhglue 块\zhglue ：}\textbf{M4/M5} \zh{｜} \zh{周\zhglue 次\zhglue ：}\textbf{W5--W6\zh{（}10-12\sym{→}10-25\zh{）}} \zh{｜} \zh{算\zhglue 力\zhglue ：}\textbf{T2 Kaggle 2\sym{×}T4\zh{，\zhglue 计\zhglue 划} 12--21 GPU\sym{·}\zh{小\zhglue 时}} \zh{唐\zhglue 杰\zhglue 作\zhglue 业\zhglue 对\zhglue 应\zhglue ：}\textbf{\sym{①} \zh{端\zhglue 到\zhglue 端\zhglue 训\zhglue 一\zhglue 个} 0.1B} \zh{｜} \zh{判\zhglue 分\zhglue 来\zhglue 源\zhglue ：}\textbf{\zh{自\zhglue 建}}\zh{（}loss/bpb \zh{达\zhglue 标} + \zh{采\zhglue 样\zhglue 可\zhglue 读\zhglue 性\zhglue ）}+ \passthrough{\lstinline!nanochat/core\_eval.py!} \zh{的} CORE \zh{手\zhglue 写\zhglue ：\zhglue 复\zhglue 用\zhglue 课\zhglue 题}01/02 \zh{的\zhglue 产\zhglue 物\zhglue （}tokenizer + \zh{模\zhglue 型} + \zh{优\zhglue 化\zhglue 器\zhglue ）\zhglue ，\zhglue 本\zhglue 课\zhglue 题\zhglue 新\zhglue 增}\textbf{\zh{数\zhglue 据\zhglue 管\zhglue 线\zhglue 与\zhglue 训\zhglue 练\zhglue 编\zhglue 排\zhglue 的\zhglue 判\zhglue 断}}
\end{quote}

\begin{center}\rule{0.5\linewidth}{0.5pt}\end{center}

\subsection{\zh{一\zhglue 、\zhglue 研\zhglue 究\zhglue 背\zhglue 景}}\label{ux4e00ux7814ux7a76ux80ccux666f}

``0.1B'' \zh{是\zhglue 唐\zhglue 杰\zhglue 作\zhglue 业\zhglue 里\zhglue 的\zhglue 数\zhglue 字\zhglue ，\zhglue 也\zhglue 是}\textbf{\zh{单\zhglue 卡}/\zh{双\zhglue 卡\zhglue 免\zhglue 费\zhglue 算\zhglue 力\zhglue 下\zhglue 的\zhglue 上\zhglue 限}}\zh{。\zhglue 参\zhglue 考\zhglue 点\zhglue ：}nanochat \zh{的} d12 \zh{约\zhglue 等\zhglue 于} GPT-1 \zh{量\zhglue 级} \zh{（\zhglue 官\zhglue 方\zhglue 讨\zhglue 论\zhglue 区\zhglue 原\zhglue 话\zhglue ：}``d12 \zh{是\zhglue 我\zhglue 最\zhglue 喜\zhglue 欢\zhglue 的\zhglue 模\zhglue 型\zhglue ，\zhglue 它\zhglue 是} GPT-1 \zh{的\zhglue 大\zhglue 小\zhglue ，\zhglue 约} 6 \zh{分\zhglue 钟\zhglue 训\zhglue 完}''------\zh{那\zhglue 是\zhglue 在} 8\sym{×}H100 \zh{上\zhglue ）\zhglue 。} \textbf{\zh{本\zhglue 课\zhglue 题\zhglue 的\zhglue 核\zhglue 心\zhglue 不\zhglue 是\zhglue 刷\zhglue 指\zhglue 标\zhglue ，\zhglue 而\zhglue 是\zhglue 把}''\zh{训\zhglue 练\zhglue 一\zhglue 次\zhglue 完\zhglue 整\zhglue 模\zhglue 型}''\zh{的\zhglue 每\zhglue 个\zhglue 环\zhglue 节\zhglue 亲\zhglue 自\zhglue 走\zhglue 通\zhglue 并\zhglue 记\zhglue 账\zhglue 。}}

\subsection{\zh{二\zhglue 、\zhglue 研\zhglue 究\zhglue 问\zhglue 题\zhglue （\zhglue 一\zhglue 句\zhglue 话\zhglue ）}}\label{ux4e8cux7814ux7a76ux95eeux9898ux4e00ux53e5ux8bdd}

\begin{quote}
\textbf{\zh{在} 2\sym{×}T4\zh{、}20 GPU\sym{·}\zh{小\zhglue 时\zhglue 的\zhglue 预\zhglue 算\zhglue 内\zhglue ，\zhglue 我\zhglue 能\zhglue 不\zhglue 能\zhglue 从\zhglue 零\zhglue 训\zhglue 出\zhglue 一\zhglue 个}''\zh{能\zhglue 生\zhglue 成\zhglue 连\zhglue 贯\zhglue 英\zhglue 文\zhglue 、\zhglue 并\zhglue 在} CORE \zh{上\zhglue 有\zhglue 非\zhglue 零\zhglue 分}''\zh{的} 0.1B \zh{模\zhglue 型\zhglue ？\zhglue 实\zhglue 测\zhglue 成\zhglue 本\zhglue 与\zhglue 估\zhglue 算\zhglue 差\zhglue 多\zhglue 少\zhglue ？}}
\end{quote}

\subsection{\zh{三\zhglue 、\zhglue 攻\zhglue 关\zhglue 线\zhglue 索}}\label{ux4e09ux653bux5173ux7ebfux7d22}

\begin{enumerate}
\def\labelenumi{\arabic{enumi}.}
\tightlist
\item
  \textbf{\zh{数\zhglue 据}}\zh{：}TinyStories\zh{（\zhglue 全\zhglue 量\zhglue ）}+ owt-sample \zh{采\zhglue 样\zhglue 。}\textbf{\zh{先\zhglue 算\zhglue 一\zhglue 遍}}\zh{：\zhglue 词\zhglue 表\zhglue 大\zhglue 小} \sym{×} \zh{语\zhglue 料\zhglue 字\zhglue 节} \sym{→} token \zh{数} \sym{→} \zh{需\zhglue 要\zhglue 的\zhglue 步\zhglue 数}
\item
  \textbf{\zh{规\zhglue 模\zhglue 选\zhglue 择}}\zh{：}nanochat \zh{用\zhglue 单\zhglue 一\zhglue 旋\zhglue 钮} \passthrough{\lstinline!--depth!}\zh{；\zhglue 在} T4 \zh{上} 16GB \zh{显\zhglue 存\zhglue 能\zhglue 放\zhglue 多\zhglue 深\zhglue ？\zhglue （\zhglue 用} M0 \zh{的\zhglue 显\zhglue 存\zhglue 公\zhglue 式\zhglue 先\zhglue 估\zhglue ）}
\item
  \textbf{\zh{精\zhglue 度}}\zh{：}\textbf{T4 \zh{是} SM75\zh{，\zhglue 没\zhglue 有} bf16} \sym{→} \zh{必\zhglue 须} fp16 + GradScaler\zh{（\zhglue 或\zhglue 用} nanochat \zh{的} \passthrough{\lstinline!NANOCHAT\_DTYPE=float16!}\zh{）}
\item
  \textbf{\zh{断\zhglue 点\zhglue 续\zhglue 训}}\zh{：}Kaggle \zh{会\zhglue 话} 9--12h \zh{会\zhglue 断} \sym{→} checkpoint \zh{必\zhglue 须\zhglue 写} \passthrough{\lstinline!/kaggle/working!}\zh{（}20GB \zh{持\zhglue 久\zhglue ）}
\item
  \textbf{\zh{超\zhglue 时\zhglue 自\zhglue 停}}\zh{：\zhglue 脚\zhglue 本\zhglue 内\zhglue 置} wall-clock \zh{上\zhglue 限\zhglue ，\zhglue 防\zhglue 止\zhglue 忘\zhglue 记\zhglue 关\zhglue 会\zhglue 话}
\item
  \textbf{\zh{指\zhglue 标}}\zh{：\zhglue 训\zhglue 练\zhglue 看} bpb\zh{（\zhglue 词\zhglue 表\zhglue 无\zhglue 关\zhglue ）\zhglue ，\zhglue 最\zhglue 终\zhglue 跑} CORE\zh{（}22 \zh{数\zhglue 据\zhglue 集\zhglue 合\zhglue 成\zhglue ，}\passthrough{\lstinline!core\_eval.py!} \zh{单\zhglue 文\zhglue 件\zhglue ）}
\end{enumerate}

\subsection{\zh{四\zhglue 、\zhglue 里\zhglue 程\zhglue 碑}}\label{ux56dbux91ccux7a0bux7891}

\begin{itemize}
\tightlist
\item[$\square$]
  M1 \textbf{\zh{实\zhglue 验\zhglue 卡\zhglue 获\zhglue 批}}\zh{：\zhglue 写\zhglue 清\zhglue 假\zhglue 设}/\zh{命\zhglue 令}/\zh{预\zhglue 算\zhglue 上\zhglue 限}/\zh{停\zhglue 止\zhglue 条\zhglue 件}/\zh{成\zhglue 功\zhglue 判\zhglue 据\zhglue （\zhglue 见\zhglue 《\zhglue 算\zhglue 力\zhglue 与\zhglue 成\zhglue 本\zhglue 预\zhglue 算}.md\zh{》}\sym{§}\zh{五\zhglue ）}
\item[$\square$]
  M2 \textbf{\zh{会\zhglue 话\zhglue 预\zhglue 检}}\zh{：}\passthrough{\lstinline!nvidia-smi!} \zh{确\zhglue 认} 2\sym{×}T4\zh{（}\textbf{\zh{若\zhglue 是} P100 \zh{立\zhglue 即\zhglue 重\zhglue 开\zhglue 会\zhglue 话}}\zh{）}
\item[$\square$]
  M3 \textbf{\zh{冒\zhglue 烟}}\zh{：}\passthrough{\lstinline!--depth=4!}\zh{、}100 \zh{步\zhglue ，\zhglue 确\zhglue 认\zhglue 能\zhglue 跑\zhglue 、\zhglue 能\zhglue 存\zhglue 、\zhglue 能\zhglue 续\zhglue 训\zhglue （}\sym{≤}10 \zh{分\zhglue 钟\zhglue ）}
\item[$\square$]
  M4 \zh{首\zhglue 次\zhglue 正\zhglue 式\zhglue 训\zhglue 练\zhglue ：\zhglue 双} T4 + DDP\zh{，\zhglue 跑\zhglue 到\zhglue 预\zhglue 算\zhglue 上\zhglue 限\zhglue 或} loss \zh{平\zhglue 台}
\item[$\square$]
  M5 \zh{续\zhglue 训\zhglue 一\zhglue 次\zhglue （\zhglue 验\zhglue 证} checkpoint \zh{真\zhglue 的\zhglue 可\zhglue 用} ------ \zh{这\zhglue 一\zhglue 条\zhglue 最\zhglue 容\zhglue 易\zhglue 翻\zhglue 车\zhglue ）}
\item[$\square$]
  M6 \zh{评\zhglue 测\zhglue ：}bpb + CORE + \textbf{\zh{采\zhglue 样} 20 \zh{条\zhglue 样\zhglue 本}}\zh{（\zhglue 人\zhglue 工\zhglue 判\zhglue 断}''\zh{是\zhglue 否\zhglue 连\zhglue 贯}''\zh{）}
\item[$\square$]
  M7 \textbf{\zh{成\zhglue 本\zhglue 账\zhglue 单}}\zh{：\zhglue 实\zhglue 测\zhglue 耗\zhglue 时} vs \zh{估\zhglue 算\zhglue （}M0 \zh{的} 6ND + T4 \zh{折\zhglue 算\zhglue ）\zhglue ，\zhglue 逐\zhglue 项\zhglue 解\zhglue 释\zhglue 偏\zhglue 差}
\item[$\square$]
  M8 \zh{一\zhglue 页\zhglue 报\zhglue 告} + \zh{知\zhglue 识\zhglue 卡\zhglue （\zhglue 含}''\zh{我\zhglue 最\zhglue 意\zhglue 外\zhglue 的\zhglue 三\zhglue 个\zhglue 现\zhglue 象}''\zh{）}
\end{itemize}

\subsection{\zh{五\zhglue 、\zhglue 交\zhglue 付\zhglue 物}}\label{ux4e94ux4ea4ux4ed8ux7269}

{\def\LTcaptype{none} % do not increment counter
\begin{longtable}[]{@{}
  >{\raggedright\arraybackslash}p{(\linewidth - 4\tabcolsep) * \real{0.3333}}
  >{\raggedright\arraybackslash}p{(\linewidth - 4\tabcolsep) * \real{0.3333}}
  >{\raggedright\arraybackslash}p{(\linewidth - 4\tabcolsep) * \real{0.3333}}@{}}
\toprule\noalign{}
\begin{minipage}[b]{\linewidth}\raggedright
\zh{交\zhglue 付\zhglue 物}
\end{minipage} & \begin{minipage}[b]{\linewidth}\raggedright
\zh{位\zhglue 置}
\end{minipage} & \begin{minipage}[b]{\linewidth}\raggedright
\zh{验\zhglue 收}
\end{minipage} \\
\midrule\noalign{}
\endhead
\bottomrule\noalign{}
\endlastfoot
\zh{实\zhglue 验\zhglue 卡} & \passthrough{\codezh{{\zh{证}}{\zh{据}}/{\zh{实}}{\zh{验}}{\zh{卡}}-NN.md}} & \zh{经\zhglue 用\zhglue 户\zhglue 确\zhglue 认} \\
run.yaml + \zh{日\zhglue 志} & \passthrough{\codezh{{\zh{证}}{\zh{据}}/}} & \zh{含} commit\zh{、\zhglue 配\zhglue 置\zhglue 、\zhglue 随\zhglue 机\zhglue 种\zhglue 子\zhglue 、\zhglue 耗\zhglue 时} \\
\zh{曲\zhglue 线} & \passthrough{\codezh{{\zh{证}}{\zh{据}}/}} & loss/bpb vs FLOPs \zh{与} vs \zh{时\zhglue 间} \zh{两\zhglue 张} \\
\zh{采\zhglue 样\zhglue 样\zhglue 本} & \passthrough{\codezh{{\zh{证}}{\zh{据}}/{\zh{样}}{\zh{本}}.md}} & \sym{≥}20 \zh{条\zhglue ，\zhglue 标\zhglue 注}''\zh{通\zhglue 顺}/\zh{不\zhglue 通\zhglue 顺}'' \\
\textbf{\zh{成\zhglue 本\zhglue 账\zhglue 单}} & \passthrough{\codezh{{\zh{报}}{\zh{告}}.md}} & \zh{实\zhglue 测} vs \zh{估\zhglue 算\zhglue 偏\zhglue 差} \sym{≤} 2\sym{×} \zh{或\zhglue 给\zhglue 出\zhglue 解\zhglue 释} \\
checkpoint & Kaggle \zh{持\zhglue 久\zhglue 盘} & \zh{不\zhglue 进\zhglue 程\zhglue （\zhglue 体\zhglue 积\zhglue 原\zhglue 因\zhglue ）\zhglue ；\zhglue 只\zhglue 在} \passthrough{\codezh{{\zh{证}}{\zh{据}}/}} \zh{留\zhglue 哈\zhglue 希\zhglue 与\zhglue 路\zhglue 径} \\
\end{longtable}
}

\subsection{\zh{六\zhglue 、\zhglue 讲\zhglue 课\zhglue 锚\zhglue 点}}\label{ux516dux8bb2ux8bfeux951aux70b9}

\begin{itemize}
\tightlist
\item
  \textbf{\zh{讲\zhglue 义}}\zh{：}\passthrough{\codezh{M4-{\zh{训}}{\zh{练}}.md}}\zh{（}W3 \zh{展\zhglue 开\zhglue ）}
\item
  \textbf{\zh{对\zhglue 标}}\zh{：}\passthrough{\lstinline!nanochat/README.md!} \zh{与} \passthrough{\lstinline!runs/speedrun.sh!}\zh{（\zhglue 读\zhglue 它\zhglue 怎\zhglue 么\zhglue 组\zhglue 织\zhglue 全\zhglue 流\zhglue 程\zhglue ）}\sym{·} CS336 lecture 05/13--14
\item
  \textbf{\zh{搭\zhglue 桥}}\zh{：\zhglue 课\zhglue 题}03 \zh{的} scaling \zh{结\zhglue 论\zhglue 在\zhglue 这\zhglue 里\zhglue 第\zhglue 一\zhglue 次\zhglue 被}''\zh{真\zhglue 跑}''\zh{检\zhglue 验}
\end{itemize}

\subsection{\zh{七\zhglue 、\zhglue 代\zhglue 码\zhglue 级\zhglue 提\zhglue 问\zhglue 示\zhglue 例}}\label{ux4e03ux4ee3ux7801ux7ea7ux63d0ux95eeux793aux4f8b}

\begin{enumerate}
\def\labelenumi{\arabic{enumi}.}
\tightlist
\item
  \zh{双} T4 \zh{跑} d12\zh{，\zhglue 为\zhglue 何}\textbf{\zh{不\zhglue 是}} 2 \zh{倍\zhglue 加\zhglue 速\zhglue ？\zhglue 先\zhglue 预\zhglue 测\zhglue 扩\zhglue 展\zhglue 效\zhglue 率\zhglue （}50\%\zh{？}80\%\zh{？\zhglue ）\zhglue ，\zhglue 再\zhglue 测\zhglue 。}
\item
  T4 \zh{上} fp16 \zh{与} fp32\zh{，\zhglue 谁\zhglue 更\zhglue 快\zhglue ？\zhglue 谁\zhglue 更\zhglue 稳\zhglue ？}GradScaler \zh{在\zhglue 解\zhglue 决\zhglue 什\zhglue 么\zhglue ？}
\item
  \zh{你\zhglue 的} loss \zh{在\zhglue 第} 2000 \zh{步\zhglue 开\zhglue 始\zhglue 上\zhglue 升\zhglue 。\zhglue 按} M0 \sym{§}\zh{五} \zh{的\zhglue 顺\zhglue 序\zhglue ，\zhglue 你\zhglue 的\zhglue 前\zhglue 三\zhglue 个\zhglue 动\zhglue 作\zhglue 是\zhglue 什\zhglue 么\zhglue ？}
\item
  \zh{若\zhglue 把\zhglue 语\zhglue 料\zhglue 从} TinyStories \zh{换\zhglue 成\zhglue 中\zhglue 文\zhglue 语\zhglue 料\zhglue （\zhglue 词\zhglue 表\zhglue 不\zhglue 变\zhglue ）\zhglue ，}loss \zh{会\zhglue 怎\zhglue 么\zhglue 变\zhglue ？}\textbf{bpb \zh{呢\zhglue ？}}
\end{enumerate}

\subsection{\zh{八\zhglue 、\zhglue 风\zhglue 险\zhglue 与\zhglue 提\zhglue 醒}}\label{ux516bux98ceux9669ux4e0eux63d0ux9192}

\begin{itemize}
\tightlist
\item
  \emo{🚨} \textbf{P100 \zh{陷\zhglue 阱}}\zh{：\zhglue 本\zhglue 课\zhglue 题\zhglue 如\zhglue 果\zhglue 拿\zhglue 到} P100\zh{，\zhglue 训\zhglue 练\zhglue 慢} 3--5 \zh{倍\zhglue 且\zhglue 后\zhglue 续\zhglue 课\zhglue 题}05 \zh{直\zhglue 接\zhglue 不\zhglue 可\zhglue 用} \sym{→} \textbf{\zh{宁\zhglue 可\zhglue 等} T4 \zh{时\zhglue 段}}\zh{。}
\item
  \emo{🚨} \textbf{\zh{配\zhglue 额}}\zh{：\zhglue 本\zhglue 课\zhglue 题\zhglue 占\zhglue 全\zhglue 学\zhglue 期\zhglue 约} 1/5 \zh{的\zhglue 配\zhglue 额\zhglue ，\zhglue 实\zhglue 验\zhglue 卡\zhglue 必\zhglue 须\zhglue 写\zhglue 实\zhglue ；}\textbf{\zh{冒\zhglue 烟\zhglue 不\zhglue 通\zhglue 过\zhglue 不\zhglue 许\zhglue 开\zhglue 正\zhglue 式\zhglue 训\zhglue 练}}\zh{。}
\item
  \sym{⚠️} 0.1B \zh{在} T4 \zh{上\zhglue 大\zhglue 概\zhglue 率\zhglue 需\zhglue 要}''\zh{降} batch + \zh{梯\zhglue 度\zhglue 累\zhglue 积}''\zh{，\zhglue 这\zhglue 会\zhglue 改\zhglue 变\zhglue 有\zhglue 效} batch size \sym{→} \textbf{\zh{要\zhglue 写\zhglue 进} run.yaml}\zh{。}
\item
  \sym{⚠️} Kaggle \zh{会\zhglue 话\zhglue 到\zhglue 期\zhglue 会\zhglue 杀\zhglue 进\zhglue 程\zhglue ：}\textbf{\zh{每} 500 \zh{步\zhglue 存\zhglue 一\zhglue 次}}\zh{是\zhglue 硬\zhglue 要\zhglue 求\zhglue 。}
\end{itemize}

\end{document}